% fig-d4-four-shadows.tex --- D4, the four cubic shadows billed by type.
%
% Defines \ClebschFourShadows.
%
% A redesign of fig:four-cubic-shadows in
% papers/clebsch-passages/sections/05-golden-operator.tex.  It follows the
% paragraph "What the four descriptions do and do not assert" that precedes
% thm:operator-shadows, which separates the four agreements into three
% registers:
%   * middle-exterior diagonal and commutator Pfaffian are ONE identity,
%     holding for every symmetric matrix;
%   * the cross-golden determinant is a reformulation once C^2 = 5I;
%   * the triangle holonomy joining them is NOT formal --- it fails for a
%     general symmetric zero-diagonal matrix, and its holding characterizes
%     the conference class.  That link carries the emphasis.
%
% The classical targets of the original figure (the six outer Joubert
% coordinates, the Segre cubic, the Segre--Igusa polar map, and the diagonal
% Clebsch cubic) are deliberately not repeated here; this figure is about
% the billing of the four agreements.

\providecommand{\ClebschFourShadows}{%
\begin{tikzpicture}[
  note/.style={font=\scriptsize,align=left,anchor=north,inner sep=0pt},
  every node/.style={outer sep=0pt},
]

\node[csfcentre,text width=34mm] (z) at (0,0.3)
  {the marked cubic\\[1pt] \(Z_T=J_T\)};

\node[csfnode,line width=1pt,fill=black!8,text width=34mm] (tri) at (-5.7,0.3)
  {triangle holonomy\\[1pt] \(\sum_S\epsilon_T(S)\,x_S\)};

\node[csfnode,text width=34mm] (wedge) at (5.2,1.65)
  {middle exterior power\\[1pt] \(\tfrac14\sum_S(K_T)_{SS}\,x_S\)};
\node[csfnode,text width=34mm] (pf) at (5.2,-0.85)
  {commutator Pfaffian\\[1pt] \(\tfrac14\operatorname{Pf}[D_x,C_T]\)};

\node[csfnode,text width=44mm] (det) at (0,-2.75)
  {oriented cross-golden determinant\\[1pt]
   \(\pm10\sqrt5\,\det B_T(x)\)};

% Register 1: the one identification that carries content.
\draw[csfstrong] (tri) -- node[font=\scriptsize,above=3pt] {\textbf{not formal}} (z);

% Register 2: one universal identity written twice.
\draw[csfbiarr] (wedge) -- node[csflab,right=1pt,text width=17mm]
  {\textbf{universal}} (pf);
\draw[thin] (wedge.west) -- ($(z.north east)+(0,-0.12)$);
\draw[thin] (pf.west) -- ($(z.south east)+(0,0.12)$);

% Register 3: a reformulation, not a further identity.
\draw[csfweak] (z) -- node[csflab,right=1pt] {\textbf{reformulation}} (det);

\node[note,text width=40mm] (ntri) at (-5.7,-1.05)
  {For a general symmetric zero-diagonal matrix the complementary minor
   \(\det A[S^{c},S]\) and the triangle product \(a_{ij}a_{jk}a_{ki}\) are
   unrelated: one uses the entries joining \(S\) to its complement, the
   other only the entries inside \(S\).  That they are proportional here is
   a property of the conference class, and it is the one identification
   among the four that could fail.};

\node[note,text width=40mm] (nuniv) at (5.2,-2.05)
  {Expanding \(\operatorname{Pf}[D_x,A]\) over perfect matchings makes the
   coefficient of \(x_S\) the signed complementary minor
   \(\det A[S^{c},S]=(*\!\bigwedge^{3}A)_{SS}\), for \emph{every} symmetric
   \(A\).  The two boxes are one identity written twice, and their
   agreement records the orientation convention rather than any property of
   \(C_T\).};

\node[note,text width=54mm] (ndet) at (0,-3.7)
  {Once \(C_T^{2}=5I\) is known this is a reformulation of the golden
   splitting.  Its content is that \(B_T(x)\) is defined without choosing
   bases, so its determinant is a map of determinant lines and the
   transported orientation fixes the sign.};

\end{tikzpicture}%
}
